\section{Conclusion}
\label{sec:conclusion}
This survey compared methods for wheeled autonomy through the navigation assumptions affected by degradation, the information or actions changed by each response, and the evidence for their effectiveness. Selective updates, independent motion sources, model adaptation, and map revision depend on different operating conditions, namely surviving geometry, credible supervision, valid correspondences, and consistent correlations.

The reviewed results also reveal conditional trade-offs, since more observable exploration routes can cost additional travel and computation, risk-constrained planning trades conservatism against progress, and recovery helps when its predicted failure matches the obstacle to task execution. Semantic ranking remains limited by the available candidates and by the capabilities of the executing chassis, whereas compact or asynchronous communication preserves only the information encoded in the shared representation. These dependencies explain why improved component scores alone do not establish more reliable navigation.

The central requirement for further evaluation is to connect degradation onset, an internal trigger, and an executed response within the same mission, as the four interfaces of Section~\ref{sec:challenges} illustrate. The coded sample shows how unevenly this connection is made, since \DegTaskShare\% of the technical publications in the purposively selected sample report a task outcome under a named degradation condition in the same experiment, from \LocSameShare\% of localization and mapping studies to \PlnSameShare\% of planning and control studies. The share is \WhlShare\% among the \GroupWHL{} publications evaluated on target wheeled robots or their data, while only \PhysSameCount{} publications obtain the combination on a physical robot. Comparisons that retain the trial unit, reference dependence, timing, and failure conditions can establish when an integrated response improves both task performance and safety. The heterogeneous, purposively selected evidence supports these mechanism-level conclusions, whereas broader performance rankings require matched protocols and more uniform coverage.
